\documentclass[10pt,a4paper]{amsart}
\usepackage[margin=19mm]{geometry}
\usepackage{amsmath,amssymb,mathtools}
\usepackage[scr=rsfs,cal=euler]{mathalfa}
\usepackage{xfrac}
\usepackage[unicode,hidelinks,bookmarks=false]{hyperref}
\newcommand{\ie}{i.e.\ }
\newcommand{\cf}{cf.\ }
\newcommand{\resp}{respectively\ }
\newcommand{\Subsection}{\subsection}
\providecommand{\nobreakdash}{\nobreak\hbox{-}}
\setlength{\emergencystretch}{2em}
\multlinegap0pt
\usepackage{bm}
\usepackage{bbm}
\usepackage{dsfont}
\usepackage{xspace}
\usepackage{cancel}
\usepackage{mathtools}
\usepackage{amsthm}
\makeatletter
\@ifundefined{theorem}{\newtheorem{theorem}{Theorem}}{}
\@ifundefined{lemma}{\newtheorem{lemma}[theorem]{Lemma}}{}
\@ifundefined{proposition}{\newtheorem{proposition}[theorem]{Proposition}}{}
\makeatother
\title[Linear programming for finite-horizon vector-valued Markov d]{Linear programming for finite-horizon vector-valued Markov decision processes}
\author{Anas Mifrani, Dominikus Noll}
\date{Source: arXiv:2502.13697v3 (CC BY 4.0)}
\begin{document}
\maketitle

\noindent\emph{Source and licence.} Linear programming for finite-horizon vector-valued Markov decision processes. Version arXiv:2502.13697v3 (CC BY 4.0), distributed under the Creative Commons Attribution 4.0 International licence, as recorded in the arXiv licence metadata for this exact version and shown on the abstract page https://arxiv.org/abs/2502.13697v3. Full rights notice: licenses/arxiv-2502.13697-CC-BY-4.0.txt.\par\smallskip

\noindent\textit{Editorial scope.} Counted unit: the Proposition whose source label is \texttt{prop2} in the article (source0.tex lines 781--786, source label \texttt{prop2}), delivered together with the article's own complete original proof of it (source0.tex lines 788--878, one \texttt{proof} environment of 91 lines). No mathematics is written, repaired or re-derived: statement and proof are the article's own text line for line. Required context retained uncounted: lem2 (lines 716--725); q\_formula (lines 721--724). Every definition, display and label of those lines is retained unchanged. Omitted from this package and not redistributed: 1--715, 725--780, 787--787, 879--2082. Nothing inside a retained excerpt is omitted or altered: every retained body is delivered exactly as the article's lines are written, so no omission receipt of any kind accompanies this package. Citations: the retained excerpts carry no citation command at all, so no bibliography entry of the article is transcribed and no reference is redistributed. Reference adaptation: none was needed and none was made; every reference inside the delivered unit resolves inside it, so no unresolved reference or citation is produced. Printed slips kept literal: no printed word, symbol, hypothesis or number of the counted statement or of its proof was altered. Layout and class: the article's own class is \texttt{article} with packages including graphicx, hyperref, indentfirst, csquotes, etoolbox, subcaption, mathrsfs, blkarray, amsmath, amssymb, amsthm, enumerate, algorithm, algorithmic; the wrapper is the lane's pinned \texttt{amsart} editorial wrapper at 10pt on A4, which restates the theorem-like environments the retained text needs and adds the article's own macro definitions that the retained text needs (none), with extra wrapper packages (bm, bbm, dsfont, xspace, cancel, mathtools). The delivered numbering is the unit's own. Assets: no retained span contains a figure environment, a graphics-inclusion command, a PDF-page inclusion, a table or a TikZ picture; no image file is copied into this package, the package declares no asset dependency and no image file is redistributed with it. Licence: version arXiv:2502.13697v3 of this article is distributed under the Creative Commons Attribution 4.0 International licence, as recorded in the arXiv licence metadata for this exact version and shown on the abstract page https://arxiv.org/abs/2502.13697v3. A full rights notice is delivered at licenses/arxiv-2502.13697-CC-BY-4.0.txt. Characters: every retained excerpt is pure ASCII, so the delivered engine, XeLaTeX with its default Latin Modern text font, reports no missing character.
\begin{lemma}
\label{lem2}
Let $\pi = (\pi_1, ..., \pi_{T-1}) \in \Pi$, $t = 1, ..., T-1$, $s \in \mathbb{S}$ and $a \in \mathbb{A}_s$.
Suppose $s$ is reachable at epoch $t$ under $\pi$. Then $\sum_{a' \in \mathbb{A}_s}x_{\pi, t}(s, a') > 0$ %by Lemma \ref{lem1}(i), 
and
\begin{equation}
\label{q_formula}
q(a \,|\, s,\ \pi_t) = \frac{x_{\pi, t}(s, a)}{{\sum_{a' \in \mathbb{A}_s}x_{\pi, t}(s, a')}}.
\end{equation}
\end{lemma}

\begin{proposition}
\label{prop2}
%
Given $x \in  P$,  there exists a unique regular policy $\pi^x \in \Pi_r$ such that $x = x_{\pi^x}$. In particular, $\pi \mapsto x_\pi$ is surjective 
onto $P$.
\end{proposition}

\begin{proof}
Uniqueness is clear from Lemma \ref{lem2}. Existence is
proved by induction.

We say that a state $s$ is \textit{attainable} by the state-action frequency vector $x$
at epoch $t$ if $\sum_{a\in \mathbb{A}_s} x_t(s,a) > 0$, otherwise $s$ is unattainable by $x$ at epoch $t$. We prove
by induction on $t$ that for every $t$ there exists a regular policy $\pi^t$ such that $x_{t'}(s,a) = x_{\pi^t,t'}(s,a)$ for all $t'=1,\dots,  t$, all $s\in \mathbb{S},a\in \mathbb{A}_s$, such that in addition (\ref{q_formula}) holds for all  $s$ attainable at $t'=1,\dots,t$.

To start the induction, note that by condition (a) we have
$\sum_{a\in \mathbb{A}_j} x_1(j,a) = \alpha(j) > 0$  for every $j$, hence every $j$ is attainable
at $t=1$ and we can put
$$
q(a|s,\pi_1^1) = \frac{x_1(s,a)}{\sum_{a'\in \mathbb{A}_s} x_1(s,a')},
$$
which fixes $\pi_1^1$.
Defining $\pi_2^1,\dots,\pi_{T-1}^1$ arbitrarily gives a policy $\pi^1$ which may without loss 
be assumed regular, because if it is not regular at some $(s,t)$ with $t >1$, we can apply the regularization procedure described earlier.
We check that
$x_{\pi^1,1}(s,a) = x_1(s,a)$ for all $s$ and $a$. Now
\begin{flalign*}x_{\pi^1, 1}(s, a) &= \sum_{j \in \mathbb{S}}\alpha(j)\,\mathbb{P}^{\pi^1}(\textbf{S}_1 = s, \textbf{A}_1 = a \,|\, \textbf{S}_1 = j) \\
&= \sum_{j \in \mathbb{S}}\alpha(j)\,q(a \,|\, s,\ \pi_1^1)\mathbb{P}^{\pi^1}(\textbf{S}_1 = s \,|\, \textbf{S}_1 = j)
\\
&= \frac{x_1(s, a)}{\sum_{a' \in \mathbb{A}_s}x_1(s, a')}\sum_{j \in \mathbb{S}}\alpha(j)\,\mathbb{P}^{\pi^1}(\textbf{S}_1 = s \,|\, \textbf{S}_1 = j) \\
%&= \frac{x_1(s, a)}{\sum_{a' \in \mathbb{A}_s}x_1(s, a')}\cdot\alpha(s) \\
&= \frac{x_1(s, a)}{\sum_{a' \in \mathbb{A}_s}x_1(s, a')}\sum_{a' \in \mathbb{A}_s}x_{\pi^1, 1}(s, a') \\
&= x_1(s, a),
\end{flalign*}
for all $s$, where the final equality is due to the fact that $\sum_{a \in \mathbb{A}_s}x_{1}(s, a) = \sum_{a \in \mathbb{A}_s}x_{\pi^1, 1}(s, a)$. Our claim therefore holds for $t = 1$. 



For the induction step, let $t=2,\dots,T-1$, and suppose there exists a regular policy $\pi^{t-1}$ 
which realizes $x$ up to epoch $t-1$, that is, 
$x_{t'}(s,a) = x_{\pi^{t-1},t'}(s,a)$ for all epochs $t'=1,\dots, t-1$ and all $s,a$, such that formula (\ref{q_formula}) is satisfied
for states $s$ reachable by $\pi^{t-1}$ at some epoch $t'=1,\dots, t-1$, while
$q(1|s,\pi^{t-1}_{t'}) = 1$ for unreachable $s$ at $t'$. 
%Now using rule (b) and the induction hypothesis, we get
%$$
%\sum_{a\in \mathbb{A}_s}x_t(s,a) =
%\sum_{j\in \mathbb{S}} \sum_{a\in \mathbb{A}_j} p_{t-1}(s|j,a) x_{t-1}(j,a)
%= \sum_{j\in \mathbb{S}} \sum_{a\in \mathbb{A}_j} p_{t-1}(s|j,a) x_{\pi^{t-1},t-1}(j,a)
%= \sum_{a\in \mathbb{A}_s} x_{\pi^{t-1},t}(s,a)
%$$
%so that attainable by $x$ at $t$ agrees with reachable by $\pi^{t-1}$ at $t$. 

%{\color{red} $q^{\pi^{t-1}_{t'}}(a_1|s)$ looks clumsy}

Now we have to define a
new regular policy $\pi^t$ which realizes $x$ up to epoch $t$. We let $\pi^t_{t'} = \pi^{t-1}_{t'}$ 
for epochs $t'=1,\dots, t-1$, and define
$$
q(a|s,\pi^t_t) = \frac{x_t(s,a)}{\sum_{a'\in \mathbb{A}_s} x_t(s,a')}
$$
for those states $s$ attainable by $x$ at epoch $t$. For
states $s$ unattainable at $t$ we define
$q(1|s,\pi_t^t)=1$. We prolong $\pi_{t+1}^t,\dots,\pi_{T-1}^t$ arbitrarily so that $\pi^t$
is regular. We now have to show that $x_{t'} = x_{\pi^t,t'}$ for $t'=1,\dots, t$, and that (\ref{q_formula}) is satisfied up to epoch $t$
at reachable states $s$. Since $\pi^t$ agrees with $\pi^{t-1}$ up to $t-1$, it remains to
check $x_t = x_{\pi^t,t}$. 
%
We have the sum equality
\begin{align*}
    \sum_{a'\in \mathbb{A}_s} x_t(s,a') &= \sum_{j\in \mathbb{S}} \sum_{a'\in \mathbb{A}_j} p_{t-1} (s|j,a') x_{t-1}(j,a') 
    \qquad\qquad \mbox{ (by (b))}\\
    &= \sum_{j\in \mathbb{S}} \sum_{a'\in \mathbb{A}_j} p_{t-1} (s|j,a') x_{\pi^{t-1},t-1}(j,a')  \qquad \mbox{ (by induction hyp.)}\\
    &=\sum_{j\in \mathbb{S}} \sum_{a'\in \mathbb{A}_j} p_{t-1} (s|j,a') x_{\pi^t,t-1}(j,a') \qquad \;\,\,\,\mbox{ (by definition of $\pi^t$})\\
    &= \sum_{a'\in \mathbb{A}_s} x_{\pi^t,t}(s,a').
    \end{align*}

    Now if $s$ is not attainable by $x$ at $t$, the left hand side equals zero, which implies that the
    right hand side also equals zero, so that the
    $x_{\pi^t,t}(s,a')$ are all zero. Hence the claimed equality is satisfied in that case.
    

    Consider the case where $s$ is attainable by $x$ at epoch $t$. Then 
    \begin{align*}
    x_{\pi^t,t}(s,a) &= \sum_{j \in \mathbb{S}} \alpha(j) \mathbb P^{\pi^t}(\textbf{S}_t=s,\textbf{A}_t=a|\textbf{S}_1=j)\\
    &= \sum_{j\in \mathbb{S}} \alpha(j) q(a|s,\pi^t_t) \mathbb P^{\pi^t}(\textbf{S}_t=s|\textbf{S}_1=j)\\
    &= \frac{x_t(s,a)}{\sum_{a'\in \mathbb{A}_s} x_t(s,a')} \sum_{j\in \mathbb{S}} \alpha(j) \mathbb P^{\pi^t}(\textbf{S}_t=s|\textbf{S}_1=j)\\
    &=  \frac{x_t(s,a)}{\sum_{a'\in \mathbb{A}_s} x_t(s,a')} \sum_{a'\in \mathbb{A}_s} x_{\pi^t,t}(s,a')\\
    &= x_t(s,a)
    \end{align*}
    where the last line uses the sum equality above. Therefore $\pi^t$ is as claimed, because
    (\ref{q_formula}) is satisfied by construction. 

    Having completed the induction,
    $\pi^{T-1}$ is a regular policy which realizes $x$ up to epoch $t=T-1$.  But since constraint (c) is explicit, $\pi^{T-1}$
    gives the desired equality also at the terminal epoch $T$, so $\pi^{T-1}$ realizes $x$ up to $T$, and
    is thus the desired
    $\pi^x$ with $x = x_{\pi^x}$.
\end{proof}

\end{document}
